% !TeX root = main.tex
\clearpage
\pagenumbering{Roman}
\setcounter{page}{1}

\chapter*{ANEXOS}
\addcontentsline{toc}{chapter}{\protect\numberline{5}ANEXOS}

\section*{ANEXO A. MANUAL DE INSTALACIÓN Y REPRODUCIBILIDAD DEL SIMULADOR}
\addcontentsline{toc}{section}{ANEXO A. MANUAL DE INSTALACIÓN Y REPRODUCIBILIDAD DEL SIMULADOR}
\label{anexo:manual_simulador}

Con el fin de garantizar la reproducibilidad y el análisis independiente del simulador físico desarrollado, el código de simulación ha sido consolidado en un repositorio público e independiente de GitHub, el cual contiene únicamente los scripts funcionales de la capa física y los bancos de pruebas.

\subsection*{A.1. Enlace del Repositorio}
El código fuente limpio se encuentra alojado en la plataforma GitHub en la siguiente dirección URL:
\begin{center}
    \url{https://github.com/carlosdaniel2440/v2i-ieee802.11bb-simulation}
\end{center}

\subsection*{A.2. Comandos de Consola para Instalación y Ejecución}
Para desplegar y ejecutar las simulaciones en cualquier computadora personal con Python 3.8 o superior instalado, se deben seguir los siguientes pasos de consola:

\begin{enumerate}
    \item \textbf{Descargar y clonar el repositorio de código:}
    \begin{verbatim}
    git clone https://github.com/carlosdaniel2440/v2i-ieee802.11bb-simulation.git
    cd v2i-ieee802.11bb-simulation
    \end{verbatim}
    
    \item \textbf{Crear e inicializar el entorno virtual de ejecución (aislamiento de dependencias):}
    \begin{verbatim}
    python -m venv venv
    \end{verbatim}
    Posteriormente, active el entorno según su sistema operativo:
    \begin{itemize}
        \item En sistemas Windows: \texttt{.\\venv\\Scripts\\activate}
        \item En sistemas Linux o macOS: \texttt{source venv/bin/activate}
    \end{itemize}
    
    \item \textbf{Instalar las dependencias matemáticas requeridas:}
    \begin{verbatim}
    pip install -r requirements.txt
    \end{verbatim}
    Este comando instalará las bibliotecas numéricas \texttt{numpy}, \texttt{scipy} y \texttt{matplotlib} compatibles.
    
    \item \textbf{Ejecutar el banco de pruebas cualitativo (Paso a Paso):}
    \begin{verbatim}
    python main.py
    \end{verbatim}
    Este comando ejecutará la transmisión del primer símbolo y guardará los gráficos de validación modular detallados en el Capítulo 2 y los resultados de canal ruidoso del Capítulo 3.
    
    \item \textbf{Ejecutar las curvas estadísticas Monte Carlo (BER y Throughput):}
    \begin{verbatim}
    python test_phy.py
    python test_phy_grid.py
    \end{verbatim}
\end{enumerate}

\clearpage

\section*{ANEXO B. ARQUITECTURA E INTERFAZ DE FUNCIONES DEL MOTOR DE CAPA FÍSICA}
\addcontentsline{toc}{section}{ANEXO B. ARQUITECTURA E INTERFAZ DE FUNCIONES DEL MOTOR DE CAPA FÍSICA}
\label{anexo:estructura_funciones_simulador}

El archivo \texttt{physical\_layer.py} actúa como la biblioteca o motor unificado del simulador. Todas sus funciones han sido escritas de forma modular y reutilizable, mapeando los bloques teóricos definidos en el estándar. A continuación se detallan las firmas de estas funciones de la capa física en dos partes: el transmisor y canal (Tabla \ref{tab:api_lc_tx}) y el receptor (Tabla \ref{tab:api_lc_rx}).

\renewcommand{\thetable}{B.\arabic{table}}
\setcounter{table}{0}

\begin{table}[H]
\centering
\caption{Diccionario y Estructura de Funciones del Módulo physical\_layer.py - Transmisor y Canal}
\label{tab:api_lc_tx}
\resizebox{\textwidth}{!}{%
\begin{tabular}{|l|l|p{6cm}|}
\hline
\textbf{Función en Python} & \textbf{Argumentos Principales} & \textbf{Propósito y Operación Matemática} \\ \hline
\texttt{fec\_encoder()} & \texttt{bits} (array 1D) & Codificación convolucional de tasa $1/2$ con polinomios generadores octales $G_0=7, G_1=5$. \\ \hline
\texttt{puncture()} & \texttt{bits, pattern} & Punzado de bits periódicos según máscara para tasas $2/3, 3/4$ o $5/6$. \\ \hline
\texttt{interleave()} & \texttt{bits} & Entrelazado de bloques mediante transposición matricial de 16 columnas para mitigar ruidos en ráfaga. \\ \hline
\texttt{constellation\_mapping()} & \texttt{bits, scheme} & Mapeo Gray multinivel (BPSK a 1024-QAM) con factor de escala $K_{\text{norm}} = 1/\sqrt{\frac{2}{3}(M-1)}$. \\ \hline
\texttt{apply\_hermitian\_symmetry()} & \texttt{symbols\_in} & Mapeo de simetría conjugada en frecuencia sobre subportadoras activas (DC y Nyquist a cero) para ondas temporales reales. \\ \hline
\texttt{idft\_processing()} & \texttt{symbols\_in} & Transformada IFFT matricial con validación del residuo imaginario inferior a $10^{-12}$. \\ \hline
\texttt{insert\_guard\_interval()} & \texttt{signal, cp\_len} & Inserción de muestras de prefijo cíclico (64, 128 o 256 muestras a $80\text{ MHz}$). \\ \hline
\texttt{windowing()} & \texttt{signal, roll\_off} & Suavizado temporal de flancos mediante ventana de coseno alzado para reducir lóbulos espectrales laterales. \\ \hline
\texttt{iq\_modulation()} & \texttt{signal, f\_if, f\_s} & Conversión hacia arriba multiplicando por portadora digital a Frecuencia Intermedia de 20 MHz. \\ \hline
\texttt{tx\_optical\_front\_end()} & \texttt{signal, min\_a, max\_a} & Suma de Bias DC adaptativo y recorte de amplitud por limitación lineal no-negativa del emisor LED. \\ \hline
\texttt{optical\_wireless\_channel()} & \texttt{signal, d, theta, ...} & Propagación en espacio libre (Lambertiano 3D), Beer-Lambert atmosférico y adición de ruidos térmico y de disparo solar. \\ \hline
\texttt{rx\_optical\_front\_end()} & \texttt{signal, R} & Conversión optoelectrónica en fotodetector según responsividad constante de $0.6\text{ A/W}$. \\ \hline
\end{tabular}%
}
\end{table}

\begin{table}[H]
\centering
\caption{Diccionario y Estructura de Funciones del Módulo physical\_layer.py - Receptor}
\label{tab:api_lc_rx}
\resizebox{\textwidth}{!}{%
\begin{tabular}{|l|l|p{6cm}|}
\hline
\textbf{Función en Python} & \textbf{Argumentos Principales} & \textbf{Propósito y Operación Matemática} \\ \hline
\texttt{remove\_dc\_bias()} & \texttt{signal} & Resta del promedio aritmético de la señal recibida (bloqueo de corriente continua). \\ \hline
\texttt{iq\_demodulation()} & \texttt{signal, f\_if, f\_s} & Conversión hacia abajo a banda base y filtrado Butterworth digital de 5to orden a 15 MHz. \\ \hline
\texttt{remove\_guard\_interval()} & \texttt{signal, cp\_len} & Descarte de las muestras iniciales del prefijo cíclico por símbolo OFDM. \\ \hline
\texttt{dft\_processing()} & \texttt{signal} & Transformada FFT matricial para regresar los símbolos al dominio de la frecuencia. \\ \hline
\texttt{remove\_hermitian\_symmetry()} & \texttt{symbols} & Descarte de subportadoras de guarda y conjugadas, aislando los 234 símbolos de datos. \\ \hline
\texttt{constellation\_demapping()} & \texttt{symbols, scheme} & Demapper Gray de decisión dura evaluando distancias euclidianas mínimas en fase y cuadratura. \\ \hline
\texttt{deinterleave()} & \texttt{bits} & Reordenamiento inverso de los bits de datos mediante matriz transpuesta de 16 columnas. \\ \hline
\texttt{depuncture()} & \texttt{bits, pattern} & Inserción de borrados (\textit{erasures}) neutros en las posiciones de bits punzadas. \\ \hline
\texttt{fec\_decoder()} & \texttt{bits, pattern} & Decodificador convolucional de Viterbi adaptado para omitir distancias en posiciones de borrado. \\ \hline
\end{tabular}%
}
\end{table}

\clearpage

\section*{ANEXO C. ESTRUCTURA COMPLETA DEL REPOSITORIO DE CÓDIGO}
\addcontentsline{toc}{section}{ANEXO C. ESTRUCTURA COMPLETA DEL REPOSITORIO DE CÓDIGO}
\label{anexo:estructura_archivos_simulador}

En el siguiente esquema se ilustra la distribución de directorios y el propósito de cada componente del software subido al repositorio de GitHub:

\begin{verbatim}
v2i-ieee802.11bb-simulation/               # Carpeta raiz del proyecto de simulacion
|
+-- physical_layer.py                            # Motor y biblioteca modular de capa fisica
+-- config_entorno.py                      # Centralizacion de constantes
|
+-- main.py                                # Banco de prueba cualitativo paso a paso
+-- test_phy.py                            # Caracterizacion estadistica Monte Carlo
+-- test_phy_grid.py                       # Simulacion matricial en rejilla
|
+-- requirements.txt                       # Librerias Python de soporte (NumPy, SciPy)
+-- README.md                              # Guia e instructivo tecnico del simulador
|
+-- plots/                                 # Directorio para exportacion de figuras
\end{verbatim}
